\chapter{因果数字孪生}
\label{ch:causal}

预测器给出“蒸汽加大后中心更快到达 $\alpha_m$”。解释器还要回答这条效应经过哪几条路径、在什么厚度上会反过来损伤表面品质。因果数字孪生是带干预语义的结构方程，参数来自第 \ref{ch:est} 章的后验，方程骨架来自第 \ref{ch:physics} 章，不能从相关数据里重新发现能量守恒。

\section{结构方程}

把一个控制周期内的决策量写成
\begin{equation}
\label{eq:scm}
Z_i=f_i(\Pa_i,U_i),\qquad i=1,\ldots,n_z,
\end{equation}
$U_i$ 是外生噪声。父节点由机理指定，学习阶段只允许调整 $f_i$ 的参数与少数预先声明的可选边，不允许从零用 PC 算法在一次烹饪的短序列上“发现因果”。

\begin{table}[H]
\centering
\caption{第一版因果骨架}
\begin{tabularx}{\textwidth}{llX}
\toprule
子节点 & 父节点 & $f_i$ 的物理内容 \\
\midrule
界面蒸汽压 $p_v$ & 蒸汽质量流、排汽阀、腔温 & 腔体水分平衡 \\
表面温度 $T_1$ & $p_v$、$T_\infty$、风速、加热功率、$T_2$ & \eqref{eq:twonode} 第一式 \\
内部温度 $T_2$ & $T_1$、$R_{\mathrm{int}}$、$C_2$ & \eqref{eq:twonode} 第二式 \\
表面含水 & $p_v$、$T_1$、$k_m$ & 蒸发—冷凝 \\
冷点进度 $\alpha_m$ & $T_2$ 的历史 & \eqref{eq:F} \\
表面褐变 $\alpha_b$ & $T_1$、表面含水 & \eqref{eq:kinetics} \\
结壳 $S$ & $\alpha_p$、表面含水 & \eqref{eq:structure} \\
内部热阻 $R_{\mathrm{int}}$ & $S$、厚度 & 结构到传递的边 \\
\bottomrule
\end{tabularx}
\end{table}

$R_{\mathrm{int}}$ 既是参数也是子节点：结壳会增大它。把这条边写进结构，是为了解释“越烤越热却越烤不透”。

\section{干预与相关的区分}

观测条件 $p(\alpha_m\mid \dot m_{\mathrm{st}})$ 里，蒸汽流量可能是控制器对当前低温的反应。相关因此混有“因为没熟所以多喷蒸汽”。干预分布
\begin{equation}
p\big(\alpha_m\mid\mathrm{do}(\dot m_{\mathrm{st}}=m^\star)\big)
\end{equation}
把决定 $\dot m_{\mathrm{st}}$ 的控制策略从结构方程中删去，只保留蒸汽到 $p_v$ 的物理边。孪生的积分实现干预：把执行器轨迹替换后再积分。三个计算依次为预测、干预和反事实。

\begin{enumerate}
\item \textbf{预测。} 保持当前策略，积分 $p(\vect{z}_{k:k+H}\mid\vect{y}_{0:k})$。
\item \textbf{干预。} 将某一执行器替换为候选值，其余结构不变，重新积分。
\item \textbf{反事实。} 先用已发生的 $\vect{y}_{0:k}$ 反推外生噪声 $U$（abduction），再在同一 $U$ 下替换执行器（action），积分得到“这一块食品如果当时少喷 10 秒”的轨迹。
\end{enumerate}

反事实使用当前后验的点估计噪声，并同时给出后验抽样带。解释文本必须带区间。没有区间的“因为蒸汽过高”句子不允许显示给用户。

\section{两条路径的符号}

以蒸汽质量流 $m_{\mathrm{st}}$ 对综合品质 $Q$ 的效应为例。$Q$ 的定义在第 \ref{ch:safe} 章，这里只分解路径。

\begin{align}
\text{路径 A: }& m_{\mathrm{st}}\to p_v\to q_{\mathrm{lat}}\to T_2\to\alpha_m\to Q,\\
\text{路径 B: }& m_{\mathrm{st}}\to p_v\to\varepsilon_{l,s}\to(\alpha_b,\,m_{\mathrm{loss}})\to Q,\\
\text{路径 C: }& m_{\mathrm{st}}\to\varepsilon_{l,s}\to S\to R_{\mathrm{int}}\to T_2\to\alpha_m\to Q.
\end{align}

路径 A 在冷点落后时为正，蒸汽的潜热加快内部进度。路径 B 在表面已经湿、且后续需要烤色时为负。路径 C 是慢路径：表面变干或浸湿改写结壳与热阻，再回到冷点。短探针窗口内路径 C 的增量是 $O(\tau/\tau_S)$，可以先不计入脉冲损伤；整个加工时域内必须保留，否则“越烤越不透”没有结构位置。总效应
\begin{equation}
\label{eq:path}
\frac{\mathrm{d}Q}{\mathrm{d}m_{\mathrm{st}}}
=\underbrace{\frac{\partial Q}{\partial\alpha_m}\frac{\partial\alpha_m}{\partial m_{\mathrm{st}}}\Big|_{\!S}}_{\text{A}}
+\underbrace{\frac{\partial Q}{\partial\varepsilon_{l,s}}\frac{\partial\varepsilon_{l,s}}{\partial m_{\mathrm{st}}}}_{\text{B}}
+\underbrace{\frac{\partial Q}{\partial\alpha_m}\frac{\partial\alpha_m}{\partial R_{\mathrm{int}}}\frac{\partial R_{\mathrm{int}}}{\partial m_{\mathrm{st}}}}_{\text{C}}.
\end{equation}
厚度增大使 A 变小。结壳使 B 更负，并使 C 在烤制段为负。解释输出的是这三项的数值与区间。

解释因此是两个数，不是一个重要度分数。两条路径分别积分，偏效应带后验区间。符号在抽样中保持稳定才生成语句，否则界面只保留预测区间。

\section{混杂与不可识别边}

下列因素会同时影响蒸汽用量和最终品质，必须在结构里显式出现，否则干预被污染：

\begin{itemize}
\item 食品厚度与初始温度（用户或视觉给定，进入 $C_2$ 与初值）；
\item 开门次数（扰动 $T_\infty$ 与 $p_v$）；
\item 加热管老化带来的功率—温度增益（由 $y_P$ 与腔温响应在线估计）。
\end{itemize}

结构中不出现、也不从数据里补上的边，包括调味料到风味分子的转化，以及微生物初始数量的绝对值。$\alpha_m$ 只作为热力当量。解释语句的词表只能引用已经积分的路径。

\section{与黑箱剩余时间模型的对照实验}

仿真里并列三个预测器：

\begin{enumerate}
\item 结构因果孪生，参数为后验；
\item 以相同输入训练的梯度提升或小型神经网络，直接回归剩余时间；
\item 因果孪生的结构被打乱一条关键边（去掉 $S\to R_{\mathrm{int}}$）的消融。
\end{enumerate}

评价不只看剩余时间 RMSE。还必须报告：在 do(蒸汽 $+10\%$) 的虚拟实验上，哪个模型的预测位移与真值仿真器一致。黑箱模型可以在观测分布上 RMSE 更低，同时在干预上失败。该失败是采用 SCM 的实验理由。

干预算法、外生噪声的反推步骤和不可识别清单见附录。对照试验的判据在第 \ref{ch:sim} 章。
